\section{System Design}\label{sec:design}

\subsection{Overview}\label{sec:design:overview}

Figure~\ref{fig:arch} shows the architecture of Drone4D. Three shared models sit on the left: the World Package with its geometry
query service (D1), the 4D environment field (D2), and the energy and perception models. Every consumer on the right
queries the same models through batched interfaces: the coupled decision layer (D4), the fleet runtime that produces
the simulated ground truth (D3), and the streaming access layer that feeds the operator's browser (D5). The design has
one rule that follows from F2: \emph{a quantity that two modules need is computed by one function}. The precheck and
the return trigger call the same energy model; the dynamics, the camera and the display read the same wind and optical
depth; the decision that proves a return safe and the contact stage that detects a collision read the same height map.
Table~\ref{tab:map} maps the modules to the challenges of Sect.~\ref{sec:bg:challenges}, and
Fig.~\ref{fig:matrix} lists, for each decision, which model it queries.

\begin{figure}[t]
\centering
\resizebox{\textwidth}{!}{\begingroup\newcommand{\chtagf}{\tiny\bfseries\color{oursred}}%
\begin{tikzpicture}[font=\scriptsize,
  box/.style={execute at begin node={\hyphenpenalty=10000},draw=g2,fill=white,rounded corners=1.5pt,line width=0.4pt,align=center,inner sep=2.5pt,
              minimum height=1.1cm},
  inp/.style={box,text width=2.3cm,fill=g5,draw=g4},
  mdl/.style={box,text width=2.9cm,draw=oursred,line width=0.7pt},
  run/.style={box,text width=3.9cm,draw=oursred,line width=0.7pt},
  arr/.style={-{Stealth[length=3pt]},g1,line width=0.45pt},
  ]
  % inputs
  \node[inp] (src) at (1.25,3.2) {City point clouds, no-fly zones, borders};
  \node[inp] (wx)  at (1.25,1.2) {Weather presets and keyframes};
  \node[inp] (veh) at (1.25,-0.8) {Vehicle and sensor catalogue};
  % shared models
  \node[mdl] (d1) at (4.6,3.2) {\textbf{D1} World Package and geometry query service\\[1pt] {\chtagf C1}};
  \node[mdl] (d2) at (4.6,1.2) {\textbf{D2} 4D environment field $E(x,y,z,t)$\\[1pt] {\chtagf C2}};
  \node[mdl] (en) at (4.6,-0.8) {Energy and perception models\\[1pt] {\chtagf C2, C3}};
  \foreach \a/\b in {src/d1, wx/d2, veh/en} \draw[arr] (\a) -- (\b);
  % bus
  \fill[oursred] (6.5,-1.35) rectangle (6.64,3.75);
  \node[font=\tiny,text=oursred,align=center,anchor=north] at (6.57,-1.42) {one shared model,\\ batched queries};
  \foreach \b in {d1,d2,en} \draw[arr] (\b.east) -- (\b.east -| 6.5,0);
  % runtime side
  \node[run] (d4) at (9.25,3.2) {\textbf{D4} Coupled decisions: precheck, return, scan plan, delegation\\[1pt] {\chtagf C3, C4}};
  \node[run] (d3) at (9.25,1.2) {\textbf{D3} Fleet runtime: one clock, staged 4\,ms ticks, rotating refresh\\[1pt] {\chtagf C4}};
  \node[run] (d5) at (9.25,-0.8) {\textbf{D5} Streaming access: gateway, budgeted point-cloud client\\[1pt] {\chtagf C5}};
  \draw[arr] (6.64,3.2) -- (d4.west);
  \draw[arr] (6.64,1.2) -- (d3.west);
  \draw[arr,dashed] (6.64,-0.8) -- (d5.west);
  % runtime internal arrows
  \draw[arr,oursred] ([xshift=-10pt]d4.south) -- node[left,font=\tiny,text=oursred] {commands} ([xshift=-10pt]d3.north);
  \draw[arr] ([xshift=10pt]d3.north) -- node[right,font=\tiny] {state} ([xshift=10pt]d4.south);
  \draw[arr] (d3.south) -- node[right,font=\tiny] {snapshots, events} (d5.north);
  % client
  \node[box,text width=1.6cm,fill=g5,draw=g4] (br) at (12.75,-0.8) {Browser on an operator's device};
  \draw[arr] (d5) -- (br);
  \draw[arr,dashed] (br.north) |- node[pos=0.25,right,font=\tiny,text=g1,align=left] {operator\\commands} (d4.east);
\end{tikzpicture}\endgroup
}
\caption{Architecture. Shared models (left) are queried in batches on the simulation clock by the decision layer and
the fleet runtime (solid) and streamed to the browser (dashed). The tags C1--C5 give the challenge each module addresses}
\label{fig:arch}
\end{figure}

\begin{table}[t]
\caption{Design modules, the challenges they address and the mechanisms that address them.}
\label{tab:map}
\footnotesize
\begin{tabular*}{\textwidth}{@{\extracolsep\fill}llp{6.3cm}@{}}
\toprule
Module & Challenge & Key mechanism \\
\midrule
D1 World Package, geometry service & C1 & Versioned package, max-pyramids, exact and bounded corridor queries \\
D2 4D environment field & C2 & One analytic field for wind, visibility and optics; golden vectors for every client \\
D3 Fleet runtime & C4 & One 250\,Hz clock, staged sub-rates, structure-of-arrays kernels, rotating return refresh \\
D4 Coupled decision layer & C3, C4 & Shared-energy return trigger, wind-rated precheck, level-driven scan planner \\
D5 Streaming access & C5 & Headless server, budgeted point-cloud selection with a cascade controller \\
\botrule
\end{tabular*}
\end{table}

\begin{figure}[t]
\centering
\begingroup
\newcommand{\cellq}[2]{\node[circle,fill=oursred,inner sep=1.5pt] at (#1) {};
  \node[font=\tiny,text=ink,align=center,anchor=north,text width=1.7cm,
        execute at begin node={\hyphenpenalty=10000}] at ($(#1)+(0,-0.12)$) {#2};}
\newcommand{\celln}[1]{\node[circle,draw=g3,inner sep=1.4pt] at (#1) {};}
\begin{tikzpicture}[font=\scriptsize,x=1.9cm,y=0.92cm]
  % header
  \foreach \c/\t/\d in {1/Geometry/D1, 2/Wind/D2, 3/Visibility/D2, 4/Energy/shared, 5/Perception/shared}
    \node[font=\scriptsize\bfseries,align=center] at (\c,0.75) {\t\\[-1pt]{\tiny\mdseries\color{g2}\d}};
  % rows
  \foreach \r/\t in {0/Launch precheck, 1/Return altitude, 2/Return trigger, 3/{Scan height, spacing},
                     4/Fleet size, 5/Task award}
    {\node[anchor=east,align=right,font=\scriptsize] at (0.45,-\r) {\t};
     \draw[g5,line width=0.4pt] (0.48,-\r-0.5) -- (5.55,-\r-0.5);}
  \draw[g2,line width=0.5pt] (0.48,0.4) -- (5.55,0.4);
  % row 0: precheck
  \cellq{1,0}{cruise above the route's roofs}
  \cellq{2,0}{route wind vs rating}
  \celln{3,0}
  \cellq{4,0}{out + station + return $\le$ reserve}
  \celln{5,0}
  % row 1: return altitude
  \cellq{1,-1}{highest roof on the line, detour}
  \cellq{2,-1}{headwind of candidates}
  \celln{3,-1}
  \cellq{4,-1}{climb vs detour energy}
  \celln{5,-1}
  % row 2: return trigger
  \cellq{1,-2}{climb to return altitude}
  \cellq{2,-2}{wind at return altitude}
  \celln{3,-2}
  \cellq{4,-2}{same model as precheck}
  \celln{5,-2}
  % row 3: scan
  \cellq{1,-3}{obstacle cells, line of sight}
  \celln{2,-3}
  \cellq{3,-3}{optical depth caps height}
  \celln{4,-3}
  \cellq{5,-3}{pixels for D/R/I}
  % row 4: fleet size
  \celln{1,-4}
  \cellq{2,-4}{sortie endurance in wind}
  \cellq{3,-4}{lower height, more strips}
  \cellq{4,-4}{per-sortie energy}
  \celln{5,-4}
  % row 5: award
  \cellq{1,-5}{line of sight to the task}
  \cellq{2,-5}{weather risk term}
  \cellq{3,-5}{expected detection}
  \cellq{4,-5}{precheck of the bid}
  \cellq{5,-5}{sensor capability}
\end{tikzpicture}
\endgroup

\caption{Decision-coupling matrix. A filled dot means the decision queries that model, with the query it makes; an
open dot means it does not. Every filled dot in a column is the same function call}
\label{fig:matrix}
\end{figure}

\subsection{World Package and geometry query service}\label{sec:design:world}

\paragraph{Package.} A city enters the system as a point cloud plus semantic layers (borders, no-fly zones). An
offline pipeline normalises units, axes and origin, derives a 10\,m terrain model by morphological opening, classifies
points by height above ground, tiles the visual cloud into an octree with a 12-byte quantised point encoding, and
derives a 2\,m digital surface model (DSM). The result is published atomically as a World Package whose content version is
the hash of its files. Six real cities from UrbanScene3D~\cite{lin2022urbanscene3d} and one procedural city are
built this way; a city occupies 138--192\,MB.

\paragraph{Derived geometry.} At load time the service derives, and caches under a key that combines the content
version with the derivation parameters, a clearance height map $H = \mathrm{maxfilter}_{5\times5}(\mathrm{DSM}) + 10$\,m,
max-pooled pyramids of the surface and of the height map, a dilated height-map pyramid for corridor queries, and an
8\,m raster of no-fly zones. Because $H$ is dilated and raised, a path that clears $H$ also clears the inflated
surface model; this monotonicity lets cheap coarse checks certify a path against the model and leaves fine checks for the few paths that need
them.

\paragraph{Queries.} Decisions use a small set of vectorised queries: surface and terrain height, the maximum of $H$
along a segment (exact cell traversal or a bounded pyramid sample with a 100\,m tolerance), batched line of sight (up
to 64 pairs per call), coarse path checks that classify each segment as proven safe, uncertain or violating against
both $H$ and the zones, fine path validation at 20 samples per metre, and a safe-transit profile that climbs to
$\max_{\text{line}} H + 5$\,m. The same height map drives the return altitude in D4 and the contact stage in D3, so
the decision and the ground truth cannot disagree about what a building is.

\subsection{4D environment field}\label{sec:design:env}

\paragraph{One analytic field.} The environment is a function $E(x,y,z,t)$ evaluated on demand rather than a stored
grid. Wind is composed as
\begin{equation}
\mathbf{w}(x,y,z,t) = s(t)\, f\big(z_{\mathrm{agl}}\big)\, \mathbf{e}(\theta(t)) + \textstyle\sum_k \mathbf{G}_k(x,y,z,t)
 + g(z)\,\mathbf{T}(x,y,z,t),
\end{equation}
a reference speed $s$ and direction $\theta$ shaped by a logarithmic profile
$f(z) = \ln\!\big((z-d)/z_0\big)/\ln\!\big((z_{\mathrm{ref}}-d)/z_0\big)$ with the city's roughness length $z_0$ and
displacement height $d$, plus discrete gusts $\mathbf{G}_k$ and turbulence $\mathbf{T}$ from a periodic $64^3$ box and a
Dryden filter bank~\cite{beal1993digital}; height above ground comes from the World Package's terrain model. The
atmosphere carries the meteorological optical range, the extinction coefficient, precipitation and fog-top height, and
the standard-atmosphere state. Slant optical depth $\tau(\mathbf{p}_0,\mathbf{p}_1,t)$ integrates an exponential haze
profile, a fog slab and precipitation below the cloud base, and gives the transmittance $e^{-\tau}$ that perception uses.
Twelve presets (clear to blizzard and sandstorm) set 21 scalar parameters; changes are applied as keyframes of at most
2\,KB that interpolate over time, so a weather change is a few bytes on the wire.

\paragraph{One field for every consumer.} The field is queried in batches of up to 256 points with a time argument.
The fleet runtime evaluates it at 50\,Hz for every drone and writes wind and air density into the dynamics; the
camera model evaluates optical depth per target; the decision layer evaluates it along candidate routes; and the
browser evaluates the same equations in TypeScript for the display. A shared set of golden vectors (profiles, gusts,
turbulence, keyframe evaluation, optical depth) is checked by both the Python and the TypeScript test suites, so the
display evaluates the same field the drones fly in (Sect.~\ref{sec:eval:consistency}); how faithfully the rendered
effects convey that field to an operator is a separate question we do not study here.

\subsection{Fleet runtime on one clock}\label{sec:design:runtime}

\paragraph{One clock, staged sub-rates.} All drones share one 250\,Hz simulation clock (4\,ms ticks). Work is split
into stages that each declare a period in ticks and a phase: the environment and sensors run at 50\,Hz, the
dynamics and contact at 125\,Hz, inter-drone collision at 25\,Hz, energy integration at 10\,Hz and return-plan refresh
at 5\,Hz. Phases are staggered to spread the expensive stages over different ticks. Every stage declares a CPU budget in
core-seconds per simulated second; the runtime refuses to build a pipeline whose budgets sum to more than 0.40 at
1,000 drones and real time, and it reports the measured cost of each stage once per second.

\paragraph{Structure of arrays.} The state of up to 1,024 drones lives in contiguous arrays (position, velocity,
attitude, battery, safety state), and every stage is a vectorised kernel over all drones or over a subset. The flight
controller is a PX4-like cascaded position and attitude controller; the power drawn is
$P = P_{h}\,(T/T_{h})^{1.5}$ plus a climb term, evaluated with the airspeed implied by the local wind. Contact against
the surface model, inter-drone separation with hold and resume, a flight state machine (flying, hold, return with
climb, cruise, descend and final phases, landing, emergency landing) and seeded random streams complete the ground
truth. No stage reads the wall clock, so a run is a deterministic function of its seed and its command log, and the
recorder can replay it.

\paragraph{Rotating refresh.} Return plans are the most expensive per-drone decision: each needs a corridor query
along the line home, a wind query at the return altitude and an energy integral. Instead of refreshing every drone on
every call, the 5\,Hz return stage refreshes a rotating fifth of the fleet, so each plan is at most one second old. Section~\ref{sec:eval:micro}
measures what this buys.

\subsection{Coupled decision layer}\label{sec:design:decisions}

\subsubsection{Decision awareness}\label{sec:design:awareness}

Each decision reads the shared models through a switchboard, \texttt{DecisionAwareness}, with one flag per factor:
\texttt{geometry} (corridor height map), \texttt{detour} (detour candidates), \texttt{wind} (field wind instead of
calm air), \texttt{visibility} and \texttt{los} (optical depth and line of sight in the detector quote),
\texttt{energy\_rtl} (energy-based return trigger), \texttt{precheck\_env} (field wind in the precheck) and
\texttt{shared\_energy} (return trigger on the precheck's energy model). A fixed-altitude policy reproduces the PX4
default. The flags change only what a decision \emph{sees}; the simulated ground truth is always fully coupled. This
makes every baseline and ablation of Sect.~\ref{sec:eval} a configuration of one system rather than a different
code base.

\subsubsection{Launch precheck}

Before a sortie the precheck builds the route the planner will fly, take-off, a climb to the cruise altitude
$z_c = \max(z_{\mathrm{op}}, z_{\mathrm{pad}} + 25, \max_{\text{route}} H + 5)$, the transit, the station time and the
return, and integrates the shared energy model along it with the field wind. A sortie launches only if
\begin{equation}
\mathrm{soc}_0 - \frac{E_{\mathrm{out}} + E_{\mathrm{station}} + E_{\mathrm{back}}}{E_{\mathrm{use}}} \geq r
\quad\text{and}\quad \max_{\text{route}} \lVert \mathbf{w}_{xy} \rVert \leq w_{\mathrm{rating}},
\end{equation}
with usable energy $E_{\mathrm{use}}$, reserve $r = 0.20$ and the vehicle's wind rating $w_{\mathrm{rating}}$
(13.8\,m/s for the P600 profile). The second condition encodes F2: wind matters mainly through feasibility.

\paragraph{Space-time precheck.} The environment field's keyframe defines the weather for the whole of its transition,
so a forecast front is part of the shared model. The space-time variant evaluates each leg at the time it will be
flown: the transit at launch time, the station and the return at their predicted start times $t_{\mathrm{back}} =
t_{\mathrm{out}} + t_{\mathrm{station}}$, with wind sampled along each leg over its duration. If the requested station
time would put the return into wind beyond the rating or below the reserve, the precheck shortens the station time in
15\,s steps, down to 30\,s, and declines the sortie only if even that fails. It costs one batched field query per
candidate station time.

\subsubsection{Return planning}

\paragraph{Altitude and path.} The return altitude is
$z_{\mathrm{rtl}} = \max\big(z_{\mathrm{now}},\, z_{\mathrm{home}} + 30,\, \max_{[\mathbf{p},\mathbf{h}]} H + 5\big)$:
the PX4 rule plus the corridor term that F1 shows to be necessary. When the corridor term forces a climb of more than
10\,m, the planner evaluates single-waypoint detours $\mathbf{p} + fL\mathbf{u} + kw\mathbf{n}$ to either side of the
line, each with its own corridor height and headwind, and keeps the one that shortens the return by at least 5\,s; at
most 16 drones run the detour search per refresh to bound its cost.

\paragraph{Trigger on the shared energy model.} The return starts when the remaining flight time falls below
$1.3\,t_{\mathrm{rtl}}$ or the state of charge below 7\%. A stand-alone failsafe estimates $t_{\mathrm{rtl}}$ from
distances and a cruise speed reduced by the headwind, which is the model disagreement of F2. With
\texttt{shared\_energy}, $t_{\mathrm{rtl}}$ is instead the energy-equivalent time of the planned return under the same
model the precheck uses:
\begin{equation}
t_{\mathrm{rtl}} = \frac{E_{\mathrm{climb}} + E_{\mathrm{cruise}}(\mathbf{w}) + E_{\mathrm{descent}}}{\max(\bar P,\, 0.5\,P_h)},
\end{equation}
where the cruise term integrates the power at the airspeed implied by the field wind at $z_{\mathrm{rtl}}$, the ground
speed is held unless the headwind exceeds the wind rating, and $\bar P$ is the current average power. Both decisions
now ask the same question of the same function, so the trigger fires when the energy, not a second model, runs out.

\subsubsection{Perception-driven scan planning}

\paragraph{From level to pixels.} For a target with projected critical dimension $d_c = \sqrt{A_{\mathrm{proj}}}$
seen at range $R$ through a lens of focal length $f$ on pixels of pitch $p$, the number of pixels on target is
$N = d_c f/(R\,p)$. The targeting task performance function~\cite{vollmerhausen2004new} gives the probability of a task
as $P(N) = x^{E}/(1+x^{E})$ with $x = N/(2N_{50})$ and $E = 2.7 + 0.7x$, where the Johnson cycle criteria
$N_{50} = 1, 4, 6.4$ correspond to detection, recognition and identification~\cite{sjaardema2015history}. At
$P = 0.9$ this requires 3.5, 14.0 and 22.4 pixels. Atmospheric contrast $C = C_0\,e^{-\tau(R)}$, with $\tau$ from D2,
scales the usable resolution by $k_c = \mathrm{clip}\big((C - 0.05)/(0.25 - 0.05), 0, 1\big)$, and a blocked line of
sight sets it to zero.

\paragraph{From pixels to a fleet.} For each candidate flying height $H$ above ground the planner (i) drops the cells
whose buildings rise above $H$ minus the clearance and rejects $H$ if they exceed 5\% of the area, (ii) computes the
slant range at the strip edge $R_e = H/\cos\theta_e$ ($\theta_e \leq 15^{\circ}$) and the focal length
$f \geq N_{\mathrm{req}} R_e\, p/(d_c\, k_c)$ needed there under the current optical depth, rejecting $H$ if no lens
reaches it, (iii) derives the ground sample distance, the strip width $w = \min(W g,\, 2H\tan\theta_{\max})$ and the
spacing $(1-s)w$, and (iv) lays out boustrophedon strips~\cite{choset2000coverage} at the best sweep angle. It keeps the
height with the shortest single-drone time, then picks the smallest number of drones whose balanced partition, assigned
to depots with the Hungarian method~\cite{kuhn1955hungarian}, fits a sortie budget
$T_{\max} = \min\big(t_{\mathrm{cap}},\, 0.8\,t_{\mathrm{sortie}}(\mathbf{w}) - 2\,t_{\mathrm{transit}}\big)$, where the
sortie endurance comes from the shared energy model under the current wind. Visibility therefore propagates through the
whole plan: lower visibility lowers the height, narrows the strips, lengthens the plan and raises the number of drones,
and the planner reports this cost instead of silently losing targets (F3).

\subsubsection{Capability-based task delegation}

Tasks that need a specific sensor are delegated with the contract net protocol~\cite{smith1980contract}: the
coordinator announces the task to the drones whose capability matches, each drone returns a quote, and the task is
awarded to the best score
$U = w_c\,\hat c - w_t\,\eta/\eta_0 - w_e\,E/E_0 - w_l\,\ell - w_r\,\rho$,
where the expected confidence $\hat c$ comes from the detector quote with optical depth and line of sight, $\eta$ and
$E$ come from the same precheck as a launch, $\ell$ is the drone's load and $\rho$ a weather-risk term from the field.
Infeasible quotes score $-\infty$.

\subsection{Streaming access without a server GPU}\label{sec:design:access}

The server runs the simulation headless on CPUs and never renders. A gateway publishes compact per-drone state at up
to 125\,Hz of simulation time (clients typically subscribe at 10--20\,Hz), events, keyframes of the environment field
and a 10\,Hz time message, with credit-based back-pressure, and it reports the age of the data it serves. The browser
renders the city point cloud with WebGL2 from HTTP range requests on the octree tiles and evaluates the environment
field locally for weather effects. Its point budget is chosen by two controllers. APH (adaptive prefix with hysteresis)
selects octree nodes best-first by screen-space error under a point budget $B$, always keeps the nodes required for a
coherent view, and applies hysteresis to avoid popping. CAS (a cascade controller) adapts $B$: an inner loop adjusts
$B$ in the log domain every 250\,ms from a 24-frame window against a 33.3\,ms frame target, and an outer loop moves
between seven quality rungs. The client therefore holds its frame rate on a laptop GPU or a software renderer,
without any GPU on the server (Sect.~\ref{sec:eval:access}).
